\documentclass[../main.tex]{subfiles}
\begin{document}
\chapter{Functionals}
\section{Definitions}
\begin{remark}[Classes of Functions]
  For $\Omega \subset \R^{n}$ we define the following sets of functions:
  \begin{align*}
    C^{k}(\Omega) &= \{y: \Omega \to \R: \text{$y$ is $k$ times continuously differentiable}\} \\
    C^{\infty}(\Omega) &= \{y: \Omega \to \R: \text{$y$ is smooth}\}
  \end{align*}
\end{remark}
\begin{definition}[Functional]
  For some set of functions $\mathscr{C}$, a functional is a map $F: \mathscr{C} \to \R$, $y \mapsto F[y]$.
  i.e. $F$ maps a function to a real number.

  More generally, a functional can map several functions to a real number:
  \[
    (y_1, \ldots, y_n) \mapsto F[y_1, \ldots, y_m]
  \]
\end{definition}
An important class functionals have $\Omega = [\alpha, \beta] \subset \R$, $\mathscr{C} = C^{\infty}(\Omega)$ and:
\begin{equation}
  F[y] = \int_{\alpha}^{\beta} f(y(x), y'(x), x) \d{x} \label{genFunct}
\end{equation}
for some smooth function $f$.
We could generalise this to depend on higher derivatives of $y$, but the first derivative is usually enough for most functionals we will encounter.
\section{Variation of Functionals}
We want to find which $y \in \mathscr{C}$ extremise the functional as defined above.
\subsection{Single Variable}
Consider what happens to $F[y]$ when we varying $y$ to $y + \delta y \in \mathscr{C}$:
\begin{align*}
  \delta F[y] &\equiv F[y + \delta y] - F[y] \\
              &= \int_{\alpha}^{\beta} (f(y + \delta y, y' + \delta y', x) - f(y, y', x)) \d{x} \\
              &= \int_{\alpha}^{\beta} \left(\delta y \pderiv{f}{y} + \delta y' \pderiv{f}{y'}\right) \d{x} + \cdots \\
\end{align*}
Note that in the above, $\delta y'$ is just the derivative of $\delta y$ with respect to $x$ and $\pderiv{f}{y'}$ is just the derivative of $f$ with respect to its second argument.

Integrating the second term by parts yields:
\[
  \int_{\alpha}^{\beta} \delta y' \pderiv{f}{y'} \d{x} = \eval{\delta y \pderiv{f}{y'}}{\alpha}{\beta} - \int_{\alpha}^{\beta} \deriv{}{x} \pderiv{f}{y'} \d{x}
\]
It would be convenient if the boundary term was not there so we assume that the definition of $\mathscr{C}$ includes boundary conditions that ensure $\eval{\delta y \pderiv{f}{y'}}{\alpha}{\beta} = 0$.

Such boundary conditions include:
\begin{itemize}
  \item They all have specific values at the endpoints, say $y(\alpha) = a$ and $y(\beta) = b$ for all $y \in \mathscr{C}$.
    Since $y + \delta y \in \mathscr{C}$, we have:
    \[
      (y + \delta y)(\alpha) = a + \delta y(\alpha) = a \implies \delta y(\alpha) = 0
    \]
    and similarly $\delta y(\beta) = 0$.
  \item They all have $\pderiv{f}{y'} = 0$ at $x = \alpha, \beta$.
    For example:
    \[
      f(y') = \sqrt{1 + (y')^2} \text{ and } y'(\alpha) = y'(\beta) = 0\ \forall y \in \mathscr{C}
    \]
\end{itemize}
So we can write:
\begin{align*}
  \delta F[y] &= \int_{\alpha}^{\beta} \delta y(x) \left(\pderiv{f}{y} - \deriv{}{x}\pderiv{f}{y'}\right) \d{x} + \cdots \\
              &= \int_{\alpha}^{\beta} \delta y(x) \frac{\delta F[y]}{\delta y(x)} \d{x} + \cdots
\end{align*}
where we define the \textit{functional derivative} of $F[y]$ with respect to $y(x)$ as:
\[
  \frac{\delta F[y]}{\delta y(x)} \equiv \pderiv{f}{y} - \deriv{}{x} \pderiv{f}{y'}
\]
For a given function $y \in \mathscr{C}$, this is just a function of $x$.

We therefore have that $\delta F[y]$ vanishes to 1st order for \textbf{arbitrary} $\delta y(x)$ if and only if $y(x)$ satisfies the \textit{Euler-Lagrange equation}:
\[
  \frac{\delta F[y]}{\delta y(x)} = \pderiv{f}{y} - \deriv{}{x} \pderiv{f}{y'} = 0
\]
This is a second order ODE in $y(x)$ which upon solving yields the extremising functions of the functional $F[y]$.
\subsection{Multivariate}
We can generalise this to vector functions $\vec{y}(x) = (y_1(x), y_2(x), \ldots, y_n(x))$ with $y_i \in \mathscr{C}\ \forall i$.
We consider an analogous type of functional defined as:
\[
  F[\vec{y}] = \int_{\alpha}^{\beta} f(\vec{y}, \vec{y}', x) \d{x}
\]
then:
\[
  \delta F[\vec{y}] = \int_{\alpha}^{\beta} \left[\sum_{i = 1}^{n} \delta y_i(x) \left(\pderiv{f}{y_i} - \deriv{}{x} \pderiv{f}{y_i'}\right)\right] \d{x} + \eval{\sum_{i = 1}^{n} \delta y_i \pderiv{f}{y_i}}{\alpha}{\beta} + \cdots
\]
So if we have boundary conditions such that the boundary term at the end vanishes, $\vec{y}(x)$ extremises $F$ if and only if:
\[
  \pderiv{f}{y_i} - \deriv{}{x} \pderiv{f}{y_i'} = 0\ \forall i \in \{1, \ldots, n\}
\]
That is, each of the $y_i$ must satisfy the Euler-Lagrange equation and so obtain $n$ second order ODEs for $y_1(x), \ldots, y_n(x)$.
\section{Geodesics in the Euclidean Plane} \label{geodesicSec}
What is the shortest curve between 2 points $A$ and $B$ in euclidean space?
We can define the length of a curve $y$ in euclidean space as the functional:
\[
  L[y] = \int_{C} \sqrt{\d{x}^2 + \d{y}^2} = \int_{C}  \d{\ell}
\]
where $\d{\ell} = \sqrt{\d{x}^2 + \d{y}^2}$.

We first need to decide how to parameterise this curve;
There are two standard ways of doing this:
\subsubsection{Parametrising with $x$}
We could assume that we can parameterise the curve as a function of $x$, i.e. $y = y(x)$ for $x_A \leq x \leq x_B$ with $y(x_A) = y_A$ and $y(x_B) = y_B$.

In this case:
\[
  L[y] = \int_{x_A}^{x_B} \sqrt{1 + (y')^2} \d{x}
\]
We have $f(y, y', x) = \sqrt{1 + {y'}^2}$ and so $\pderiv{f}{y} = 0$.
Note also that our boundary conditions are appropriate since we have fixed the value of $y$ at the endpoints of the curve.

Thus to satisfy the Euler-Lagrange equations we must have:
\[
  \deriv{}{x} \pderiv{f}{y'} = 0 \implies \pderiv{f}{y'} = \frac{y'}{\sqrt{1 + (y')^2}} = k
\]
for some constant $k$.
After rearranging, we obtain that $y' = c$ for some constant $c$ and so $y = mx + c$.
We can then use the boundary conditions on $y$ to obtain the values of $m$ and $c$.

This method works however we are excluding curves where $y$ \textbf{is not a function} of $x$.
\subsubsection{Arbitrary Parameter}
We can instead use an arbitrary parameter $t \in [0, 1]$ where $t = 0$ at $A$ and $t = 1$ at $B$.

We write this as a vector $\vec{x}(t) = (x(t), y(t))$ and so the functional becomes:
\[
  L[\vec{x}] = \int_{0}^{1} \sqrt{\dot{\vec{x}}(t) \cdot \dot{\vec{x}}(t)} \d{t} = \int_{0}^{1} \sqrt{\dot{x}^2 + \dot{y}^2} \d{t}
\]
with boundary conditions $\vec{x}(0) = \vec{x}_A$ and $\vec{x}(1) = \vec{x}_B$.
In this functional we have $f(x, \dot{x}, y, \dot{y}, t) = \sqrt{\dot{x}^2 + \dot{y}^2}$ and so the Euler-Lagrange equation for $x(t)$ is:
\[
  \deriv{}{t} \pderiv{f}{\dot{x}} - \pderiv{f}{x} = \pderiv{f}{\dot{x}} = 0 \implies \pderiv{f}{\dot{x}} = \frac{\dot{x}}{\sqrt{\dot{x}^2 + \dot{y}^2}}= c
\]
and similarly for $y$ by symmetry:
\[
  \pderiv{f}{\dot{y}} = \frac{\dot{y}}{\sqrt{\dot{x}^2 + \dot{y}^2}} = s
\]
for constants $c$ and $s$.

We now observe that $c^2 + s^2 = 1$ and so we can write the constants as $c = \cos \theta$ and $s = \sin \theta$ for some fixed $\theta$.
It is now convenient to switch to using \textit{arc length} $\ell$ as our parameter.
To do this we first need to find $\deriv{\ell}{t}$:
\[
  \deriv{\ell}{t} = \sqrt{\left(\deriv{x}{t}\right)^2 + \left(\deriv{y}{t}\right)^2} = \sqrt{\dot{x}^2 + \dot{y}^2}
\]
Using the chain rule we have:
\[
  \deriv{x}{\ell} = \left(\deriv{x}{t} \middle/ \deriv{\ell}{t} \right) = \frac{\dot{x}}{\sqrt{\dot{x}^2 + \dot{y}^2}} = \cos \theta
\]
and similarly:
\[
  \deriv{y}{\ell} = \left(\deriv{y}{t} \middle/ \deriv{\ell}{t} \right) = \frac{\dot{y}}{\sqrt{\dot{x}^2 + \dot{y}^2}} = \sin \theta
\]
We can immediately integrate these with respect to $\ell$ and apply the boundary conditions to yield:
\[
  \vec{x}(\ell) = \vec{x}_{A} + \ell \begin{pmatrix}
  \cos \theta \\
  \sin \theta \\
  \end{pmatrix}
\]
which is a straight line parameterised by arc length.
\begin{remark}
  We could not have used $\ell$ initially as then the upper limit of integration would depend on the curve $C$ but our derivation of the Euler-Lagrange equations assumed that we had constant limits of integration.
\end{remark}
\section{First Integrals and Fermat's Principle}
\subsection{First Integrals}
Consider again the functional:
\[
  F[y] = \int_{\alpha}^{\beta} f(y, y', x) \d{x}
\]
If we have $\pderiv{f}{y} = 0$, then the Euler-Lagrange equation is just:
\[
  \deriv{}{x} \pderiv{f}{y'} = 0
\]
which we can immediately integrate to obtain a \textit{first integral}:
\[
  \pderiv{f}{y'} = \text{constant}
\]
This is just a first order ODE and so this case is typically considerably easier to solve as otherwise we are working with a second order, usually non-linear, ODE.
\begin{proposition}
  If $f$ has no explicit $x$-dependence, that is $\pderiv{f}{x} = 0$, then: \label{noXDep}
  \[
    f - y' \pderiv{f}{y'} = \text{constant}
  \]
  for any solution of the Euler-Lagrange equation.
\end{proposition}
\begin{remark}[Note]
  $f$ still depends \textit{implicitly} on $x$ via $y(x)$.
\end{remark}
\begin{proof}
  Using the multivariate chain rule we have:
  \begin{align*}
    \deriv{f}{x} &= \deriv{}{x}(f(y(x), y'(x), x)) \\
                 &= \pderiv{f}{x} + \pderiv{f}{y}y' + \pderiv{f}{y'} y'' \\
                 &= y'\left[\pderiv{f}{y} - \deriv{}{x}\pderiv{f}{y'}\right] + \pderiv{f}{y'}y'' + y' \deriv{}{x} \pderiv{f}{y'} \text{ since $\pderiv{f}{x} = 0$}\\
                 &= \deriv{}{x}\left(y' \pderiv{f}{y'}\right) \text{ since $y$ is a solution of the EL equation}
  \end{align*}
  We now consider:
  \[
    \deriv{}{x}\left(f - y' \pderiv{f}{y'}\right) = \deriv{f}{x} - \deriv{}{x}\left(y' \pderiv{f}{y'}\right) = 0
  \]
  and so $f - y' \deriv{f}{y'} = \text{constant}$, as required.
\end{proof}
\subsection{Fermat's Principle}
Consider light propagating in a medium, for example in a glass box.
Let $c(\vec{x})$ be the speed of light in the medium at a point $\vec{x}$ and let $c$ be the constant speed of light in the vacuum.
\begin{definition}[Refractive Index]
  We define the \textit{refractive index} as:
  \[
    n(\vec{x}) = \frac{c(\vec{x})}{c}
  \]
\end{definition}
The time $T$ taken for the light to travel from $A$ to $B$ along some path $C$ is then:
\[
  T[C] = \int_{C} \frac{1}{c(\vec{x})} \d{\ell} = \frac{1}{c} P[C]
\]
where $P[C]$ is the \textit{optical path length}:
\[
  P[C] = \int_{C} n(\vec{x}) \d{\ell}
\]
\begin{law}[Fermat's Principle]
  Light takes the path from $A$ to $B$ along a path that \textbf{locally minimises} $T$, or equivalently $P$.
\end{law}
\begin{remark}
  ``Locally'' here just means that $T$ is minimised between any two points on the path which are sufficiently close together.
\end{remark}
\begin{example}[Light Rays]
  Find the path of a light ray in the $(x, z)$-plane in a medium with refractive index:
  \[
    n(z) = \sqrt{a - bz} \text{ where } a, b > 0
  \]
  For simplicity, we assume that this path can be parameterised as $z = z(x)$ for $x \in [\alpha, \beta]$ where $x$ is at $A$ at $\alpha$ and $B$ at $\beta$.
  The functional $P$ then is:
  \[
    P[z] = \int_{\alpha}^{\beta} n(z) \underbrace{\sqrt{1 + {z'}^2} \d{x}}_{\d{\ell}} = \int_{\alpha}^{\beta} f(z, z') \d{x}
  \]
  Since there is no explicit $x$ dependence in $f$, by \cref{noXDep}, we have that:
  \begin{align*}
    k &= f - z' \pderiv{f}{z'} \\
      &= n(z)\sqrt{1 + {z'}^2} - \frac{z'^2 n(z)}{\sqrt{1 + {z'}^2}} \\
      &= \frac{n(z)}{\sqrt{1 + {z'}^2}} = \sqrt{\frac{a - bz}{1 + {z'}^2}}
  \end{align*}
  for some constant $k$.
  Rearranging, we have:
  \[
    {z'}^2 = \frac{b}{k^2}(z_0 - z) \text{ where } z_0 = \frac{a - k^2}{b}
  \]
  and so:
  \[
    z' = \pm \sqrt{\frac{b}{k^2}(z_0 - z)} \implies \frac{z'}{\sqrt{z_0 - z}} = \pm \sqrt{\frac{b}{k^2}}
  \]
  We can rewrite the left hand side of this as:
  \[
    \deriv{}{x}(- 2 \sqrt{z_0 - z}) = \pm\sqrt{\frac{b}{k^2}} \implies z = z_0 - \frac{b}{4k^2}(x - x_0)^2
  \]
  for a constant of integration $x_0$.
  We can then use the boundary conditions to determine the constants of integration $k$ and $x_0$ (and then $z_0$).
  This defines a parabola in the $(x, z)$-plane.
\end{example}
\begin{example}[Brachistochrone Curve]
  Suppose we have a bead sliding on a frictionless wire on a vertical plane.
  What shape of wire minimises the time for the bead to fall from rest at $A$ to a lower point $B$ that is horizontally displaced?

  We choose our axes so that $A$ is at the origin:
  \begin{center}
  \begin{tikzpicture}[>=stealth]
    \draw[->] (0, 0) node[left] {$A$} -- (5, 0) node[right] {$x$};
    \draw[->] (0, -3) -- (0, 1) node[above] {$y$};
    \node[right] at (2.3, -1.9) {$B$};
    \draw[very thick, gray!70] (0, 0) -- (2.3, -1.9) node[midway, above right] {\large?};
  \end{tikzpicture}
  \end{center}
  We appeal to energy conservation to find the velocity at a given $y$:
  \[
    \frac{1}{2} mv^2 + mgy = 0 \implies v = \sqrt{-2gy}
  \]
  We also make the reasonable assumption that the curve does not double back on itself and so we can use $x$ as a parameter of the curve:
  \[
    T = \int_{A}^{B} \frac{\d{\ell}}{v} = \frac{1}{\sqrt{2g}} \int_{A}^{B} \frac{\sqrt{\d{x}^2 + \d{y}^2}}{\sqrt{-y}} \implies T \propto \int_{0}^{x_B} \frac{\sqrt{1 + {y'}^2}}{\sqrt{-y}} \d{x}
  \]
  Since there is no explicit $x$ dependence, we have a first integral by \cref{noXDep} of the form:
  \[
    \text{constant} = f - y' \pderiv{f}{y'} = \frac{1}{\sqrt{(-y)(1 + {y'}^2)}} \implies (-y)(1 + {y'}^2) = 2c \geq0
  \]
  for some constant of integration $c$ where we specify $2c$ for convenience.

  If we substitute $y = -c(1 - \cos \theta)$ for $\theta(x) \geq 0$ then:
  \[
    y' = \deriv{y}{\theta} \deriv{\theta}{x} = -c \theta' \sin \theta
  \]
  and regard $x$ as a function of $\theta$, i.e. $x(\theta)$ then we obtain:
  \[
    \deriv{x}{\theta} = \pm c (1 - \cos \theta)
  \]
  and so the ODE becomes:
  \begin{align*}
    (1 - \cos \theta)(1 + (c\theta' \sin \theta)) &= 2 \\
    \iff (c \theta' \sin \theta)^2 &= \frac{2}{1 - \cos \theta} - 1 \\
                                   &= \frac{\sin \theta}{(1 - \cos \theta)^2} \\
    \iff {\theta'}^2 &= \frac{1}{c^2(1 - \cos \theta)^2}
  \end{align*}
  If we instead regard $\theta$ as a function of $x$ then we need to solve:
  \[
    \left(\deriv{x}{\theta}\right)^2 = c^2(1 - \cos \theta)^2
  \]
  Choosing the positive root so that $x > 0$, we obtain:
  \[
    \deriv{x}{\theta} = c(1 - \cos \theta) \implies x = c(\theta - \sin \theta) + k
  \]
  Since $x = 0$ at $y = 0$, we must have $x = 0$ at $\theta = 0$ by our definition of $y$ and so $k = 0$.
  We can then use the boundary condition of the point $B$ to determine $c$ and the upper bound on $\theta$.
  This a parametric description of an \textit{inverted cycloid}:
  \begin{center}
  \begin{tikzpicture}[>=stealth]
    \draw[->] (0, 0) node[left] {$A$} -- (5, 0) node[right] {$x$};
    \draw[->] (0, -3) -- (0, 1) node[above] {$y$};
    \draw[very thick, rotate=180, domain=-2.6:0, smooth, variable=\t] plot ({1 * (\t - sin(\t r))},{1*(1 - cos(\t r))});
    \node at (2.3, -1.9) {$B$};
  \end{tikzpicture}
  \end{center}
  A cycloid is the path traced out by a fixed point on the rim of a disc that rolls along a straight line at a constant speed.
\end{example}
\section{Constrained Functionals}
\subsection{Functional Lagrange Multipliers}
Suppose that we wish to extremise some functional $F[y]$ but now subject to a \textit{functional constraint} $P[y] = c$.
We can use a Lagrange multiplier similar to those in \cref{lagrangeMulti} but instead with functionals:
\[
  \Phi_{\lambda} [y] = F[y] - \lambda(P[y] - c)
\]
If we extremise with respect to $y$ (assuming the boundary term in the variation vanishes), then we obtain:
\[
  \frac{\delta F}{\delta y(x)} - \lambda \frac{\delta P}{\delta y(x)} = 0
\]
which is the Euler-Lagrange equation for this problem.

If we extremise with respect to $\lambda$, then this recovers the constraint $P[y] = c$ in the same way as it did before.
\subsection{Examples}
Determine $y(x)$ that minimises the functional:
\[
  F[y] = \int_{0}^{1} \frac{1}{2} {y'}^2 \d{x}
\]
subject to the constraint:
\[
  P[y] = \int_{0}^{1} y(x) \d{x} = 1
\]
in the class of smooth functions on $[0, 1]$ with boundary values $y(0) = y(1) = 0$.

The Lagrange multiplier for this is:
\begin{align*}
  \Phi_{\lambda}[y] &= \int_{0}^{1} \frac{1}{2} {y'}^2 \d{x} - \lambda \left(\int_{0}^{1} y \d{x} - 1\right) \\
                    &= \int_{0}^{1} \left(\frac{1}{2} {y'}^2 - \lambda y\right) \d{x} + \lambda \\
                    &= \int_{0}^{1} f \d{x} + \lambda \text{ where } f(y, y', x) = \frac{1}{2} {y'}^2 - \lambda y
\end{align*}
To extremise this, the Euler-Lagrange equations for $\Phi_{\lambda}$ tell us that:
\[
  0 = \pderiv{f}{y} - \deriv{}{x} \pderiv{f}{y'} = - \lambda - y'' \implies y'' = - \lambda
\]
\begin{remark}[Note]
  We could have directly used:
  \[
    \frac{\delta F}{\delta y(x)} - \lambda \frac{\delta P}{\delta y(x)} = - \deriv{}{x}(y'') - \lambda(1 - 0) = 0 \iff y'' = -\lambda
  \]
\end{remark}
Upon solving this we obtain:
\[
  y(x) = - \frac{1}{2} \lambda x^2 + Ax + B
\]
Applying the boundary conditions, we obtain $B = 0$ and $A = \frac{\lambda}{2}$.

To find $\lambda$, we use the constraint:
\[
  1 = \int_{0}^{1} y \d{x} = \eval{-\frac{1}{6} \lambda x^3 + \frac{1}{2}\left(\frac{1}{2}\lambda\right)x^2}{0}{1} = \frac{\lambda}{12} \implies \lambda = 12
\]
and so the solution is $y = 6x(1 - x)$.
\begin{example}[Further Examples]
  \begin{itemize}
    \item \textbf{Dido's Problem --} Suppose that we have a fence of length $L$ and a wall.
      What is the optimal shape for the fence such that the area enclosed by the wall and the fence is maximised?
      See Example Sheet 1 Q12.
    \item \textbf{Isoperimetric Problem --} Suppose that a simple closed curve has length $L$, what shape should it be to maximise the area that it encloses?
      See Example Sheet 2 Q3.
  \end{itemize}
\end{example}
\subsection{Sturm-Liouville Problem}
Consider the class of functions:
\[
  \mathscr{C} = \{y: [\alpha , \beta] \to \R: y \text{ is twice continuously differentiable and } y(\alpha) = y(\beta) = 0\}
\]
and define the functionals:
\begin{align*}
  F[y] &= \int_{\alpha}^{\beta} [\rho(x) y'(x)^2 + \sigma(x) y(x)^2] \d{x} \\
  G[y] &= \int_{\alpha}^{\beta} w(x) y(x)^2 \d{x}
\end{align*}
where $\rho(x), w(x) > 0$ for $x \in (\alpha, \beta)$ and we call $w(x)$ a \textit{weight function}.
The problem is then to minimise $F[y]$ in $\mathscr{C}$ subject to the functional constraint $G[y] = 1$.

We start by introducing a Lagrange multiplier:
\[
  \Phi_{\lambda}[y] = F[y] - \lambda(G[y] - 1)
\]
It is suitable for us to use the Euler-Lagrange equations as since the endpoints are fixed, we have appropriate boundary conditions.
Thus we have:
\begin{equation}
  \frac{\delta F[y]}{\delta y(x)} - \lambda \frac{\delta G[y]}{\delta y(x)} = 0 \tag{$\star$} \label{LGSL}
\end{equation}
We see:
\begin{align*}
  \frac{\delta F[y]}{\delta y(x)} &= \pderiv{}{y}(\rho {y'}^2 + \sigma y^2) - \deriv{}{x}\left[\pderiv{}{y'}(\rho {y'}^2 + \sigma y^2)\right] \\
                                  &= 2\sigma y - \deriv{}{x}(2 \rho y') \\
                                  &\equiv 2 \mathcal{L} y
\end{align*}
where we define the differential operator $\mathcal{L}$ as:
\[
  \mathcal{L} \equiv -\deriv{}{x}\left(\rho \deriv{}{x}\right) + \sigma
\]
We also see that:
\[
  \frac{\delta G[y]}{\delta y(x)} = \pderiv{}{y}(w y^2) - \deriv{}{x}\left[\pderiv{}{y'}(w y^2)\right] = 2wy
\]
and so \cref{LGSL} becomes:
\begin{equation}
  \mathcal{L} y = \lambda w y \tag{$\dag$} \label{SLP}
\end{equation}
With boundary conditions $y(\alpha) = y(\beta) = 0$, this is an \textit{eigenvalue problem}.
Only for discrete $\lambda$ does there exist a non-trivial solution $y(x)$, known as an \textit{eigenfunction}.
This is known as a \textit{Sturm-Liouville} problem;
We will study this class of problems further in IB Methods.

We now notice that since \cref{SLP} is linear, if $y(x)$ solves then so does $Cy(x)$ for any constant $C$.
\begin{remark}
  We can think of the constraint $G[y] = 1$ as an analogue to the restriction of only wanting unit eigenvectors of a matrix.
\end{remark}
We can also integrate the functional $F[y]$ by parts to obtain:
\begin{align*}
  F[y] &= \int_{\alpha}^{\beta} [y' \cdot \rho y' + \sigma y^2] \d{x} \\
       &= \cancelto{0}{\eval{\rho y y'}{\alpha}{\beta}} + \int_{\alpha}^{\beta} \left[\sigma y^2 - y \deriv{}{x}(\rho y')\right] \d{x} \text{ as $y(\alpha) = y(\beta) = 0$} \\
       &= \int_{\alpha}^{\beta} y \mathcal{L} y \d{x} \\
       &= \lambda \int_{\alpha}^{\beta} w y^2 \d{x} \text{ by \cref{SLP}} \\
       &= \lambda G[y] = \lambda
\end{align*}
So $F[y] = \lambda$ for eigenfunctions of the above Sturm-Liouville problem.
Thus the solution to our minimisation problem is the smallest eigenvalue $\lambda$.
\begin{remark}[Remark]
  This strategy also works with the alternative boundary conditions:
  \begin{itemize}
    \item $y'(\alpha) = y'(\beta) = 0$
    \item $y(\alpha) = y'(\beta) = 0$
    \item $y'(\alpha) = y(\beta) = 0$
  \end{itemize}
\end{remark}
Similarly to what we did in \cref{alternativeApproach}, we could also consider extremising:
\[
  \Lambda[y] \equiv \frac{F[y]}{G[y]}
\]
over $y \in \mathscr{C}$ \textbf{without any constraints}.
By setting the functional derivative of $\Lambda$ to 0 we obtain:
\[
  0 = \frac{\delta \Lambda}{\delta y(x)} = \frac{1}{G}\left(\frac{\delta F[y]}{\delta y(x)} - \frac{F}{G} \frac{\delta G[y]}{\delta y(x)}\right) = \frac{2}{G}(\mathcal{L}y - \Lambda w y)
\]
and so the minimal value of $\Lambda$, and therefore also the solution to our constrained problem, is the lowest eigenvalue of the Sturm Liouville problem.
\subsection{Geodesics and Function  Constraints}
Consider a surface $S$ in Euclidean space with equation $g(\vec{x}) = 0$.
If we are given $A, B \in S$, how do we find the shortest curve \textbf{on $S$} that connects $A$ and $B$.

Let $\vec{x}(t)$ for $t \in [0, 1]$ be a curve on $S$ with $\vec{x}(0) = \vec{x}_A$ and $\vec{x}(1) = \vec{x}_B$.
We want to extremise the \textbf{functional}:
\[
  L[\vec{x}] = \int_{0}^{1} \sqrt{\dot{\vec{x}} \cdot \dot{ \vec{x}}} \d{t}
\]
now also subject to a constraint given by a \textbf{function} $g(\vec{x}(t)) = 0$.

To do this we introduce a \textit{Lagrange Multiplier Function} $\lambda(t)$:
\[
  \Phi[\vec{x}, \lambda] = L[\vec{x}] - \int_{0}^{1} \lambda(t) g(\vec{x}(t)) \d{t}
\]
If we extremise $\Phi$ with respect to $\lambda$, then the EL equations tell us:
\begin{equation}
  0 = \frac{\delta \Phi}{\delta \lambda(t)} = - \pderiv{}{\lambda}(\lambda g) + \cancelto{0}{\deriv{}{t}\left[\pderiv{}{\lambda'}(\lambda g)\right]} = - g(\vec{x}(t)) \iff g(\vec{x}(t)) = 0 \tag{$\dag$} \label{lExtremeGeo}
\end{equation}
which is the constraint.

If we extremise with respect to $x_i(t)$, then the EL equations tell us:
\begin{align}
  0 &= \frac{\delta \Phi}{\delta x_i(t)} \nonumber \\
    &= \cancelto{0}{\pderiv{}{x_i}(\sqrt{\dot{\vec{x}} \cdot \dot{\vec{x}}})} - \deriv{}{t}\left[\pderiv{}{\dot{x}_i}(\sqrt{\dot{\vec{x}} \cdot \dot{\vec{x}}})\right] - \pderiv{}{x_i} (\lambda g) + \cancelto{0}{\pderiv{}{\dot{x}_i} (\lambda g)} \nonumber \\
    &= - \lambda \nabla_{i}g(\vec{x}(t)) - \deriv{}{t}\left[\frac{\dot{x}_i}{\sqrt{\dot{\vec{x}} \cdot \dot{\vec{x}}}}\right] \tag{$\ast$} \label{xExtremeGeo}
\end{align}
We can now eliminate $\lambda$.
First we differentiate \cref{lExtremeGeo} with respect to $t$:
\[
  0 = \deriv{}{t}(g(\vec{x}(t))) = \dot{x}_i \pderiv{g}{x_i} \implies \frac{\dot{x}_i}{\sqrt{\dot{\vec{x}} \cdot \dot{\vec{x}}}} \pderiv{g}{x_i} = 0
\]
and upon differentiating with respect to $t$ again:
\begin{align*}
  \deriv{}{t} \left[\frac{\dot{x}_i}{\sqrt{\dot{\vec{x}} \cdot \dot{\vec{x}}}}\right] \pderiv{g}{x_i} + \frac{\dot{x}_i \dot{x}_j}{\sqrt{\dot{\vec{x}} \cdot \dot{\vec{x}}}} \frac{\partial^2 g}{\partial x_j \partial x_i} = 0 \\
\end{align*}
By \cref{xExtremeGeo}, we have that:
\[
  \deriv{}{t} \left[\frac{\dot{x}_i}{\sqrt{\dot{\vec{x}} \cdot \dot{\vec{x}}}}\right] \pderiv{g}{x_i} = - \lambda \nabla_i g \nabla_i g = -\lambda |\nabla g|^2
\]
and so we can solve for $\lambda(t)$:
\[
  \lambda(t) = \frac{\dot{x}_i \dot{x}_j}{|\nabla g|^2 \sqrt{\dot{\vec{x}} \cdot \dot{\vec{x}}}} \frac{\partial^2 g}{\partial x_j \partial x_i}
\]
Finally, dividing \cref{xExtremeGeo} by $\sqrt{\dot{\vec{x}} \cdot \dot{\vec{x}}}$ and substituting $\lambda(t)$ we have:
\[
  \frac{1}{\sqrt{\dot{\vec{x}} \cdot \dot{\vec{x}}}}\deriv{}{t} \left[\frac{\dot{x}_i}{\sqrt{\dot{\vec{x}} \cdot \dot{\vec{x}}}}\right] + \frac{\dot{x}_j \dot{x}_k}{|\nabla g|^2 \dot{\vec{x}} \cdot \dot{\vec{x}}} \frac{\partial^2 g}{\partial x_j \partial x_k} \nabla_i g = 0\ \forall i
\]
or equivalently:
\begin{equation}
  \frac{1}{|\dot{\vec{x}}|} \deriv{}{t}\left[\frac{\dot{\vec{x}}}{|\dot{\vec{x}}|}\right] + \frac{\dot{\vec{x}}^{\trans} H_g \dot{\vec{x}}}{|\nabla g|^2 |\dot{\vec{x}}|^2}  \nabla g = \vec{0} \label{geodesicEqn}
\end{equation}
where $H_g$ is the Hessian of $g(\vec{x})$.
This is the equation for geodesics on a surface which is solved along with the constraint in \cref{lExtremeGeo}.

We can clean this up by converting to arc length as a parameter:
\[
  \deriv{}{\ell} = \frac{1}{|\dot{\vec{x}}|} \deriv{}{t} \implies \vec{x}' = \frac{\dot{\vec{x}}}{|\dot{\vec{x}}|}
\]
where we use $\vec{x}'$ to denote differentiation with respect to $\ell$.
Thus \cref{geodesicEqn} becomes:
\[
  \vec{x}'' + \frac{{\vec{x}'}^{\trans} H_g \vec{x}'}{|\nabla g|^2} \nabla g = \vec{0}
\]
\begin{example}
  If we set $g = z$, then $S$ is the plane $z = 0$ and so from \cref{geodesicSec}, we expect the result to be a straight line.

  We see that:
  \[
    [H_g]_{i j} = \frac{\partial^2 g}{\partial x_i \partial x_j} = 0 \implies \vec{x}'' = \vec{0}  \implies \vec{x} = \vec{c} + \ell \vec{d}
  \]
  which is a straight line.
  We can then use the boundary conditions to determine $\vec{c}$, $\vec{d}$ and the range for $\ell$.
\end{example}
\end{document}
